\documentclass[12pt]{article}

\usepackage{amsmath,amssymb,amsthm,mathrsfs}
\usepackage{enumitem}
\usepackage{geometry}
\usepackage{hyperref}
\usepackage{mathtools}

\geometry{margin=1in}
\setlength{\parskip}{0.65em}
\setlength{\parindent}{0pt}

% Theorem styles
\newtheorem{theorem}{Theorem}[section]
\newtheorem{lemma}[theorem]{Lemma}
\newtheorem{proposition}[theorem]{Proposition}
\newtheorem{corollary}[theorem]{Corollary}

\theoremstyle{definition}
\newtheorem{definition}[theorem]{Definition}
\newtheorem{example}[theorem]{Example}
\newtheorem{remark}[theorem]{Remark}

\newcommand{\norm}[1]{\left\lVert #1\right\rVert}
\newcommand{\supnorm}[1]{\left\lVert #1\right\rVert_\infty}

\title{\textbf{Real Analysis Lecture Notes}\\
\large Sequences of Functions, Uniform Convergence, Differentiation/Integration, Power Series, Abel's Theorem}
\author{}
\date{}

\begin{document}
\maketitle
\tableofcontents
\newpage

%%%%%%%%%%%%%%%%%%%%%%%%%%%%%%%%%%%%%%%%%%%%%%%%%%%%%%%%%%%%
\section*{Source}
These notes are compiled and expanded into full lectures (with complete proofs and worked solutions) from the question set in the PDF \texttt{Real-2.pdf}. 

%%%%%%%%%%%%%%%%%%%%%%%%%%%%%%%%%%%%%%%%%%%%%%%%%%%%%%%%%%%%
\section{Lecture Plan (Time Allocation)}
Assume each lecture is \textbf{90 minutes}. A suggested pacing:

\begin{enumerate}[label=\textbf{L\arabic*.}]
\item \textbf{(90 min)} Pointwise vs uniform convergence; uniform $\Rightarrow$ pointwise; key examples.
\item \textbf{(90 min)} Cauchy criterion; testing uniform convergence by supremum estimates.
\item \textbf{(90 min)} Pointwise but not uniform; when pointwise $\Rightarrow$ uniform (Dini, equicontinuity ideas).
\item \textbf{(90 min)} Uniform convergence and differentiation: theorem + proof + applications.
\item \textbf{(90 min)} Uniform convergence and integration: theorem + proof; counterexamples when uniform fails.
\item \textbf{(90 min)} Advanced examples: spikes, oscillations; boundedness of derivatives vs uniform convergence.
\item \textbf{(90 min)} Recursive sequences of functions $f_{n+1}=\sqrt{x+f_n(x)}$; uniform convergence via Dini.
\item \textbf{(90 min)} Power series: radius of convergence; transformations; Abel's theorem + applications to $\log2$ and $(\log(1+x))^2$.
\end{enumerate}

\newpage

%%%%%%%%%%%%%%%%%%%%%%%%%%%%%%%%%%%%%%%%%%%%%%%%%%%%%%%%%%%%
\section{Lecture 1: Pointwise and Uniform Convergence}

\subsection{Definitions}

\begin{definition}[Pointwise Convergence]
Let $E\subseteq \mathbb{R}$ and $f_n:E\to\mathbb{R}$. We say $f_n\to f$ \emph{pointwise} on $E$ if
for every $x\in E$ and every $\varepsilon>0$, there exists $N=N(\varepsilon,x)$ such that
\[
|f_n(x)-f(x)|<\varepsilon\quad\text{for all }n\ge N.
\]
\end{definition}

\begin{definition}[Uniform Convergence]
We say $f_n\to f$ \emph{uniformly} on $E$ if for every $\varepsilon>0$, there exists $N=N(\varepsilon)$ such that
\[
|f_n(x)-f(x)|<\varepsilon\quad\text{for all }x\in E\text{ and all }n\ge N.
\]
Equivalently,
\[
\sup_{x\in E}|f_n(x)-f(x)|\xrightarrow[n\to\infty]{}0.
\]
\end{definition}

\subsection{Uniform Convergence Implies Pointwise Convergence}

\begin{theorem}
If $f_n\to f$ uniformly on $E$, then $f_n\to f$ pointwise on $E$.
\end{theorem}

\begin{proof}
Fix $x\in E$ and $\varepsilon>0$. Uniform convergence gives an $N(\varepsilon)$ such that
$|f_n(t)-f(t)|<\varepsilon$ for all $t\in E$ and all $n\ge N(\varepsilon)$.
In particular, for this fixed $x$ we have $|f_n(x)-f(x)|<\varepsilon$ for all $n\ge N(\varepsilon)$.
Thus pointwise convergence holds at $x$. Since $x$ was arbitrary, pointwise convergence holds on $E$.
\end{proof}

\subsection{Example: A Basic Uniform Limit Question}

\begin{example}
Consider $f_n:(0,1)\to\mathbb{R}$ given by
\[
f_n(x)=\frac{n x}{n^2+1}\qquad(n\ge 1).
\]
Does it have a uniform limit? Find it and justify.
\end{example}

\begin{proof}[Solution]
For each fixed $x\in(0,1)$,
\[
f_n(x)=x\cdot \frac{n}{n^2+1}\xrightarrow[n\to\infty]{}0,
\]
so the pointwise limit is $f(x)=0$.

To test uniform convergence on $(0,1)$, estimate the sup:
\[
\sup_{x\in(0,1)}|f_n(x)-0|
=\sup_{x\in(0,1)} \frac{n x}{n^2+1}
= \frac{n}{n^2+1}\sup_{x\in(0,1)} x
= \frac{n}{n^2+1}.
\]
Since $\frac{n}{n^2+1}\to 0$, we get $\sup_{x\in(0,1)}|f_n(x)|\to 0$.
Hence $f_n\to 0$ uniformly on $(0,1)$.
\end{proof}

%%%%%%%%%%%%%%%%%%%%%%%%%%%%%%%%%%%%%%%%%%%%%%%%%%%%%%%%%%%%
\section{Lecture 2: Cauchy Criterion and Uniform Convergence Tests}

\subsection{Cauchy Criterion}

\begin{theorem}[Cauchy Criterion for Uniform Convergence]
A sequence $f_n$ converges uniformly on $E$ if and only if
for every $\varepsilon>0$ there exists $N$ such that
\[
|f_n(x)-f_m(x)|<\varepsilon\quad\text{for all }x\in E\text{ and all }n,m\ge N.
\]
Equivalently,
\[
\sup_{x\in E}|f_n(x)-f_m(x)|\xrightarrow[n,m\to\infty]{}0.
\]
\end{theorem}

\begin{proof}
($\Rightarrow$) If $f_n\to f$ uniformly then for given $\varepsilon>0$ pick $N$ so that
$\sup_{x\in E}|f_k(x)-f(x)|<\varepsilon/2$ for all $k\ge N$. Then for $n,m\ge N$,
\[
|f_n(x)-f_m(x)|\le |f_n(x)-f(x)|+|f(x)-f_m(x)|<\varepsilon.
\]
Taking sup over $x$ gives the Cauchy condition.

($\Leftarrow$) Assume the Cauchy condition. Fix $x\in E$. Then $(f_n(x))$ is a Cauchy sequence in $\mathbb{R}$,
so it converges to some $f(x)$. Define $f:E\to\mathbb{R}$ by this pointwise limit.
Now fix $\varepsilon>0$ and choose $N$ such that $\sup_{x\in E}|f_n(x)-f_m(x)|<\varepsilon$ for all $n,m\ge N$.
Let $m\to\infty$; since $f_m(x)\to f(x)$ for each fixed $x$, we get $|f_n(x)-f(x)|\le \varepsilon$ for all $x\in E$ and $n\ge N$.
Hence $\sup_{x\in E}|f_n(x)-f(x)|\le \varepsilon$, proving uniform convergence.
\end{proof}

\subsection{Test Problems (Worked)}

\begin{example}
Test convergence of $f_n(x)=\dfrac{x}{1+n x^2}$ on $\mathbb{R}$.
\end{example}

\begin{proof}[Solution]
Pointwise: for fixed $x$,
\[
f_n(x)=\frac{x}{1+n x^2}\xrightarrow[n\to\infty]{}
\begin{cases}
0,& x\neq 0,\\
0,& x=0,
\end{cases}
\]
so the pointwise limit is $0$.

Uniform: compute the supremum. For $x\neq 0$, consider $g(x)=\frac{|x|}{1+n x^2}$ for $x\ge 0$.
Differentiate for $x>0$:
\[
g'(x)=\frac{(1+n x^2)-x(2n x)}{(1+n x^2)^2}
=\frac{1+n x^2-2n x^2}{(1+n x^2)^2}
=\frac{1-n x^2}{(1+n x^2)^2}.
\]
Thus maximum occurs at $x=\frac{1}{\sqrt{n}}$ and
\[
\sup_{x\in\mathbb{R}}|f_n(x)|
= g\!\left(\frac{1}{\sqrt{n}}\right)
= \frac{\frac{1}{\sqrt{n}}}{1+n\cdot \frac{1}{n}}
= \frac{1/\sqrt{n}}{2}
= \frac{1}{2\sqrt{n}}
\xrightarrow[n\to\infty]{}0.
\]
Hence $f_n\to 0$ uniformly on $\mathbb{R}$.
\end{proof}

\begin{example}
Test convergence of $f_n(x)=\dfrac{\sin(nx)}{\sqrt{n}}$ on $[0,2\pi]$.
\end{example}

\begin{proof}[Solution]
For each fixed $x$, $|\sin(nx)|\le 1$, hence $|f_n(x)|\le 1/\sqrt{n}\to 0$.
So pointwise limit is $0$.

For uniform convergence,
\[
\sup_{x\in[0,2\pi]}|f_n(x)|\le \frac{1}{\sqrt{n}}\to 0.
\]
Thus $f_n\to 0$ uniformly on $[0,2\pi]$.
\end{proof}

\begin{example}
Test convergence of $f_n(x)=n x e^{-n x^2}$ on $[0,\infty)$ (and on $[0,1]$).
\end{example}

\begin{proof}[Solution]
Pointwise: for any fixed $x>0$, $n x^2\to\infty$ so $e^{-n x^2}\to 0$ faster than any power of $n$;
hence $n x e^{-n x^2}\to 0$. Also $f_n(0)=0$. So pointwise limit is $0$ on $[0,\infty)$.

Uniform: write $x\ge 0$ and set $t=\sqrt{n}\,x$. Then
\[
f_n(x)=n x e^{-n x^2} = \sqrt{n}\, t e^{-t^2}.
\]
Now $\phi(t)=t e^{-t^2}$ has maximum where $\phi'(t)=e^{-t^2}(1-2t^2)=0$,
so at $t=1/\sqrt{2}$, and $\max_{t\ge 0}\phi(t)=\frac{1}{\sqrt{2e}}$.
Thus
\[
\sup_{x\ge 0} f_n(x)=\sqrt{n}\cdot \frac{1}{\sqrt{2e}} \xrightarrow[n\to\infty]{}\infty.
\]
Therefore $f_n\not\to 0$ uniformly on $[0,\infty)$.
The same supremum point occurs at $x=\frac{1}{\sqrt{2n}}\le 1$ for large $n$,
so also $f_n\not\to 0$ uniformly on $[0,1]$.
\end{proof}

%%%%%%%%%%%%%%%%%%%%%%%%%%%%%%%%%%%%%%%%%%%%%%%%%%%%%%%%%%%%
\section{Lecture 3: Pointwise vs Uniform; When Pointwise $\Rightarrow$ Uniform}

\subsection{Classic Counterexample (Fully Worked)}

\begin{example}
Define $f_n:[0,1]\to\mathbb{R}$ by
\[
f_n(x)=
\begin{cases}
n x, & 0\le x\le \frac{1}{n},\\[4pt]
\frac{1}{n}, & \frac{1}{n}< x\le 1.
\end{cases}
\]
Show $f_n\to 0$ pointwise but not uniformly on $[0,1]$.
\end{example}

\begin{proof}[Solution]
Fix $x\in(0,1]$. For all large $n$ we have $x>\frac{1}{n}$, so $f_n(x)=\frac{1}{n}\to 0$.
Also $f_n(0)=0$ for all $n$. Hence $f_n(x)\to 0$ for every $x\in[0,1]$ (pointwise).

For uniform convergence, compute $\sup_{x\in[0,1]}|f_n(x)|$.
On $[0,1/n]$, $f_n(x)=nx$ achieves maximum $1$ at $x=1/n$.
Thus
\[
\sup_{x\in[0,1]}|f_n(x)| = 1\quad\text{for all }n,
\]
so $\sup_{x}|f_n(x)-0|\not\to 0$. Hence the convergence is not uniform.
\end{proof}

\subsection{When does pointwise convergence become uniform?}

\begin{theorem}[Dini's Theorem]
Let $K$ be compact, and $f_n:K\to\mathbb{R}$ continuous. Suppose $f_n\to f$ pointwise on $K$,
$f$ is continuous, and the convergence is monotone (either $f_n\uparrow f$ or $f_n\downarrow f$ pointwise).
Then $f_n\to f$ uniformly on $K$.
\end{theorem}

\begin{proof}
Assume $f_n\uparrow f$. Define $g_n=f-f_n\ge 0$. Then $g_n$ are continuous and $g_n\downarrow 0$ pointwise.
If not uniform, there exists $\varepsilon_0>0$ and points $x_n\in K$ with $g_n(x_n)\ge \varepsilon_0$.
By compactness, pass to a subsequence $x_{n_k}\to x_\ast\in K$.
Since $g_{n_k}$ is decreasing in $k$ pointwise, for any fixed $m$ and $k$ large enough with $n_k\ge m$,
\[
g_m(x_{n_k})\ge g_{n_k}(x_{n_k})\ge \varepsilon_0.
\]
Let $k\to\infty$; continuity of $g_m$ gives $g_m(x_\ast)\ge \varepsilon_0$ for all $m$.
But $g_m(x_\ast)\downarrow 0$ pointwise, contradiction. Hence uniform.
\end{proof}

%%%%%%%%%%%%%%%%%%%%%%%%%%%%%%%%%%%%%%%%%%%%%%%%%%%%%%%%%%%%
\section{Lecture 4: Uniform Convergence and Differentiation}

\subsection{Main Theorem + Complete Proof}

\begin{theorem}[Uniform Convergence and Differentiation]
Let $f_n$ be differentiable on $[a,b]$. Assume:
\begin{enumerate}[label=(\arabic*)]
\item There exists $x_0\in[a,b]$ such that $(f_n(x_0))$ converges in $\mathbb{R}$.
\item $f_n'\to g$ uniformly on $[a,b]$ for some function $g$.
\end{enumerate}
Then $f_n$ converges uniformly on $[a,b]$ to a differentiable function $f$ and $f'=g$ on $[a,b]$.
\end{theorem}

\begin{proof}
Fix $x\in[a,b]$. By the Fundamental Theorem of Calculus,
\[
f_n(x)-f_n(x_0)=\int_{x_0}^{x} f_n'(t)\,dt.
\]
Let $m,n$ be large. Then
\[
f_n(x)-f_m(x) = (f_n(x_0)-f_m(x_0)) + \int_{x_0}^{x}\big(f_n'(t)-f_m'(t)\big)\,dt.
\]
Taking absolute values and using $| \int_{x_0}^x h|\le |b-a|\sup_{[a,b]}|h|$,
\[
|f_n(x)-f_m(x)| \le |f_n(x_0)-f_m(x_0)| + (b-a)\sup_{t\in[a,b]}|f_n'(t)-f_m'(t)|.
\]
Now $(f_n(x_0))$ is Cauchy, and $f_n'$ is uniformly Cauchy (since it converges uniformly). Hence the RHS
can be made $<\varepsilon$ for all $x\in[a,b]$ by choosing $m,n$ large.
Therefore $f_n$ is uniformly Cauchy on $[a,b]$, so it converges uniformly to some $f$.

Next, since $f_n'\to g$ uniformly and $f_n(x)-f_n(x_0)=\int_{x_0}^x f_n'(t)\,dt$, take $n\to\infty$.
Uniform convergence of derivatives implies
\[
\int_{x_0}^x f_n'(t)\,dt \to \int_{x_0}^x g(t)\,dt
\]
and uniform convergence of $f_n$ implies $f_n(x)-f_n(x_0)\to f(x)-f(x_0)$.
Thus
\[
f(x)-f(x_0) = \int_{x_0}^x g(t)\,dt.
\]
Hence $f$ is differentiable and $f'(x)=g(x)$ for all $x$ (by FTC).
\end{proof}

\subsection{Application 1: $f_n(x)=\frac{\cos(nx)}{n}$}

\begin{example}
Let $f_n(x)=\frac{\cos(nx)}{n}$ on $\mathbb{R}$. Show $f_n\to 0$ uniformly.
Find where $f_n'$ converges. Construct a modified example where $f_n\to 0$ uniformly but $\{f_n'\}$ is unbounded.
\end{example}

\begin{proof}[Solution]
Uniform convergence:
\[
\sup_{x\in\mathbb{R}}|f_n(x)| \le \frac{1}{n}\to 0,
\]
so $f_n\to 0$ uniformly on $\mathbb{R}$.

Derivative:
\[
f_n'(x)= -\sin(nx).
\]
The sequence $\sin(nx)$ converges only at points $x$ where $\sin(nx)$ is constant in $n$.
If $x=k\pi$ for integer $k$, then $\sin(nk\pi)=0$ for all $n$, so $f_n'(x)\to 0$.
If $x\notin \pi\mathbb{Z}$, then $\sin(nx)$ oscillates and does not converge (one can show it has at least two limit points).
Thus $f_n'$ converges exactly on $\pi\mathbb{Z}$ (to $0$).

Now modify to make derivatives unbounded while $f_n\to 0$ uniformly:
take
\[
g_n(x)=\frac{\sin(n^2 x)}{n}.
\]
Then $\sup_x |g_n(x)|\le 1/n\to 0$ so $g_n\to 0$ uniformly.
But
\[
g_n'(x)= n\cos(n^2 x),
\]
so $\sup_x |g_n'(x)| = n\to \infty$: derivatives are unbounded.
\end{proof}

\subsection{Application 2: $f_n(x)=x-\frac{x^n}{n}$ (Correct domain discussion)}

This family behaves differently depending on the domain. On all of $\mathbb{R}$ it does not even converge pointwise for $|x|>1$,
since $\frac{x^n}{n}$ diverges when $|x|>1$.

\begin{example}
On a compact interval $[-c,c]$ with $0<c<1$, apply the theorem to
\[
f_n(x)=x-\frac{x^n}{n}.
\]
\end{example}

\begin{proof}[Solution]
On $[-c,c]$,
\[
\sup_{|x|\le c}\left|\frac{x^n}{n}\right| \le \frac{c^n}{n}\to 0,
\]
so $f_n\to f$ uniformly where $f(x)=x$.

Differentiate:
\[
f_n'(x)=1-x^{n-1}.
\]
For $|x|\le c<1$,
\[
\sup_{|x|\le c}|f_n'(x)-1|=\sup_{|x|\le c}|x^{n-1}|\le c^{n-1}\to 0,
\]
so $f_n'\to 1$ uniformly.
Also $f_n(0)=0$ converges. Thus the differentiation theorem applies and yields $f'=1$,
consistent with $f(x)=x$.
\end{proof}

%%%%%%%%%%%%%%%%%%%%%%%%%%%%%%%%%%%%%%%%%%%%%%%%%%%%%%%%%%%%
\section{Lecture 5: Uniform Convergence and Integration}

\subsection{Theorem + Complete Proof}

\begin{theorem}[Uniform Convergence and Integration]
Let $f_n:[a,b]\to\mathbb{R}$ be Riemann integrable for each $n$. If $f_n\to f$ uniformly on $[a,b]$,
then $f$ is Riemann integrable and
\[
\int_a^b f(x)\,dx = \lim_{n\to\infty}\int_a^b f_n(x)\,dx.
\]
\end{theorem}

\begin{proof}
Uniform convergence means $\sup_{[a,b]}|f_n-f|\to 0$.
Then
\[
\left|\int_a^b f_n - \int_a^b f\right|
\le \int_a^b |f_n-f|
\le (b-a)\sup_{[a,b]}|f_n-f|
\xrightarrow[n\to\infty]{}0,
\]
so $\int f_n\to \int f$.
Also uniform limit of integrable functions is integrable (a standard result; can be proved using upper/lower sums).
\end{proof}

\subsection{Spikes: Why Uniform Matters (Full Computations)}

\begin{example}
Define on $[0,1]$:
\[
f_n(x)=
\begin{cases}
n(1-nx), & 0<x<\frac{1}{n},\\
0, & \text{otherwise}.
\end{cases}
\]
Compute $\int_0^1 f_n$ and compare with the pointwise limit.
\end{example}

\begin{proof}[Solution]
Pointwise, for fixed $x>0$ and $n>\frac{1}{x}$, we have $x>\frac{1}{n}$, so $f_n(x)=0$.
Also $f_n(0)=0$. Hence $f_n(x)\to 0$ pointwise on $[0,1]$.

But
\[
\int_0^1 f_n(x)\,dx = \int_0^{1/n} n(1-nx)\,dx
= \int_0^{1/n} (n-n^2 x)\,dx
= \left[nx-\frac{n^2x^2}{2}\right]_{0}^{1/n}
=1-\frac12=\frac12.
\]
So $\int_0^1 f_n = 1/2$ for all $n$, while $\int_0^1 0 = 0$.
This shows \emph{without uniform convergence} we may fail to interchange limit and integral.

Indeed, $\sup_{[0,1]} f_n = n$ (unbounded), hence no uniform convergence.
\end{proof}

\begin{example}
Let $f_n(x)=n x e^{-n x^2}$ on $[0,1]$. Compute $\int_0^1 f_n$ and compare with pointwise limit.
\end{example}

\begin{proof}[Solution]
We already proved $f_n\to 0$ pointwise but not uniformly on $[0,1]$.

Compute integral: substitute $u= n x^2$ so $du=2n x\,dx$ and $n x\,dx=\frac12 du$:
\[
\int_0^1 n x e^{-n x^2}\,dx = \frac12 \int_{0}^{n} e^{-u}\,du
= \frac12(1-e^{-n}) \xrightarrow[n\to\infty]{}\frac12.
\]
So again $\lim \int f_n = 1/2 \ne 0 = \int \lim f_n$.
Failure is explained by lack of uniform convergence.
\end{proof}

%%%%%%%%%%%%%%%%%%%%%%%%%%%%%%%%%%%%%%%%%%%%%%%%%%%%%%%%%%%%
\section{Lecture 6: Cauchy Criterion Example $\tan^{-1}(nx)$}

\begin{example}
Using the Cauchy criterion, show that $f_n(x)=\tan^{-1}(nx)$ is uniformly convergent on $[a,b]$ for $a>0$,
but not uniformly convergent on $[0,b]$.
\end{example}

\begin{proof}[Solution]
First note pointwise behavior: for $x>0$, $nx\to\infty$ so $\tan^{-1}(nx)\to \pi/2$.
At $x=0$, $f_n(0)=0$. Hence the pointwise limit on $[0,b]$ is
\[
f(x)=\begin{cases}
0,& x=0,\\
\pi/2,& x\in(0,b].
\end{cases}
\]
On $[a,b]$ with $a>0$, the pointwise limit is the constant $\pi/2$.

\textbf{Uniform on $[a,b]$, $a>0$:}
For $x\in[a,b]$, set $y=nx\ge na$. Use the estimate (monotonicity of $\tan^{-1}$):
\[
0\le \frac{\pi}{2}-\tan^{-1}(y)=\tan^{-1}\!\left(\frac{1}{y}\right)\le \frac{1}{y}
\]
(using $\tan^{-1}(u)\le u$ for $u\ge 0$).
Thus for $x\in[a,b]$,
\[
0\le \frac{\pi}{2}-f_n(x)=\frac{\pi}{2}-\tan^{-1}(nx)\le \frac{1}{nx}\le \frac{1}{na}.
\]
Hence
\[
\sup_{x\in[a,b]}\left|f_n(x)-\frac{\pi}{2}\right|\le \frac{1}{na}\to 0,
\]
so uniform convergence holds on $[a,b]$.

\textbf{Not uniform on $[0,b]$:}
If it were uniform on $[0,b]$, then the uniform limit would be continuous if each $f_n$ is continuous
(uniform limit of continuous functions is continuous).
But the pointwise limit $f$ has a jump discontinuity at $0$.
Therefore convergence cannot be uniform on $[0,b]$.
\end{proof}

%%%%%%%%%%%%%%%%%%%%%%%%%%%%%%%%%%%%%%%%%%%%%%%%%%%%%%%%%%%%
\section{Lecture 7: Recursive Sequence $f_{n+1}(x)=\sqrt{x+f_n(x)}$}

\begin{example}
Let $x\ge 0$ and define $f_1(x)=\sqrt{x}$, $f_{n+1}(x)=\sqrt{x+f_n(x)}$.
Discuss uniform convergence on (i) $[a,b]$ with $0<a<b$ and (ii) $[0,1]$.
\end{example}

\begin{proof}[Solution]
\textbf{Step 1: Monotone increasing in $n$.}
Since $f_1(x)\ge 0$ for $x\ge 0$, we have $f_2(x)=\sqrt{x+f_1(x)}\ge \sqrt{x}=f_1(x)$.
Assume $f_n(x)\ge f_{n-1}(x)\ge 0$. Then
\[
f_{n+1}(x)=\sqrt{x+f_n(x)} \ge \sqrt{x+f_{n-1}(x)}=f_n(x),
\]
so by induction $f_n(x)$ increases in $n$ for each fixed $x\ge 0$.

\textbf{Step 2: Uniform boundedness on compact intervals.}
Fix $x\in[0,1]$. We claim $f_n(x)\le 2$ for all $n$.
Indeed, if $f_n(x)\le 2$, then $f_{n+1}(x)=\sqrt{x+f_n(x)}\le \sqrt{1+2}=\sqrt{3}<2$.
Also $f_1(x)=\sqrt{x}\le 1\le 2$. Hence by induction $f_n(x)\le 2$ on $[0,1]$.

Similarly on $[a,b]$ one gets a uniform bound.

\textbf{Step 3: Pointwise limit and formula.}
Since $f_n(x)$ is increasing and bounded above, it converges pointwise to some $f(x)$.
Passing to the limit in $f_{n+1}(x)=\sqrt{x+f_n(x)}$ gives
\[
f(x)=\sqrt{x+f(x)} \quad\Rightarrow\quad f(x)^2-f(x)-x=0.
\]
Since $f(x)\ge 0$, we choose the positive root:
\[
f(x)=\frac{1+\sqrt{1+4x}}{2}.
\]
This is continuous on $[0,\infty)$.

\textbf{Step 4: Uniform convergence via Dini.}
Each $f_n$ is continuous (composition of continuous functions).
On $[a,b]$ and on $[0,1]$ (both compact), we have:
\begin{itemize}
\item $f_n$ continuous,
\item $f_n\uparrow f$ pointwise,
\item $f$ continuous.
\end{itemize}
By Dini's theorem, $f_n\to f$ uniformly on $[a,b]$ (for $0<a<b$) and also uniformly on $[0,1]$.
\end{proof}

%%%%%%%%%%%%%%%%%%%%%%%%%%%%%%%%%%%%%%%%%%%%%%%%%%%%%%%%%%%%
\section{Lecture 8: Power Series, Radius of Convergence, Abel's Theorem}

\subsection{Radius of Convergence}

\begin{definition}[Radius of Convergence]
For a power series $\sum_{n=0}^\infty a_n x^n$, define
\[
\alpha:=\limsup_{n\to\infty} |a_n|^{1/n}.
\]
Then the radius of convergence is
\[
R:=\begin{cases}
\infty,& \alpha=0,\\
\frac{1}{\alpha},& 0<\alpha<\infty,\\
0,& \alpha=\infty.
\end{cases}
\]
The series converges absolutely for $|x|<R$ and diverges for $|x|>R$.
\end{definition}

\begin{remark}
\textbf{Everywhere convergent} means $R=\infty$. \textbf{Nowhere convergent} (except possibly at $0$) means $R=0$.
Examples:
\[
\sum_{n=0}^\infty \frac{x^n}{n!} \text{ has }R=\infty,\qquad
\sum_{n=0}^\infty n!x^n \text{ has }R=0.
\]
\end{remark}

\subsection{Transforming the Radius}

\begin{theorem}
If $\sum a_n x^n$ has radius of convergence $R$, then $\sum a_n x^{kn}$ has radius of convergence $R^{1/k}$
for each positive integer $k$.
\end{theorem}

\begin{proof}
Let $\alpha=\limsup |a_n|^{1/n}$, so $R=1/\alpha$ (with conventions).
Consider $\sum a_n x^{kn}$, which is $\sum b_n x^n$ where $b_{kn}=a_n$ and $b_m=0$ otherwise.
Alternatively, use root test directly:
\[
\limsup_{n\to\infty} |a_n|^{1/n} |x|^k = \alpha |x|^k.
\]
Convergence occurs when $\alpha |x|^k<1$, i.e. $|x|<\alpha^{-1/k}=R^{1/k}$.
\end{proof}

\subsection{Computing Radii (Worked)}

\begin{example}
Find $R$ for $\displaystyle \sum_{n=1}^\infty \frac{(n-1)!}{n^n} x^n$.
\end{example}

\begin{proof}[Solution]
Use $\alpha=\limsup |a_n|^{1/n}$ with $a_n=\frac{(n-1)!}{n^n}$.
Using Stirling-type growth $(n!)^{1/n}\sim n/e$, we get
\[
((n-1)!)^{1/n}\sim \frac{n}{e}.
\]
Hence
\[
|a_n|^{1/n} = \frac{((n-1)!)^{1/n}}{n}
\to \frac{1}{e}.
\]
So $\alpha=1/e$ and $R=1/\alpha=e$.
\end{proof}

\begin{example}
Find $R$ for $\displaystyle \sum_{n=1}^\infty \frac{(n-1)^n}{n^n} x^n$.
\end{example}

\begin{proof}[Solution]
Here $a_n=\left(\frac{n-1}{n}\right)^n = \left(1-\frac1n\right)^n$.
Then
\[
|a_n|^{1/n} = 1-\frac1n \to 1.
\]
So $\alpha=1$ and $R=1$.
\end{proof}

\subsection{Abel's Theorem (Statements) and Applications}

\begin{theorem}[Abel's Theorem: First Form (Continuity at Endpoint)]
Let $\sum_{n=0}^\infty a_n x^n$ have radius $R=1$. If $\sum_{n=0}^\infty a_n$ converges, then
\[
\lim_{x\uparrow 1}\sum_{n=0}^\infty a_n x^n = \sum_{n=0}^\infty a_n.
\]
\end{theorem}

\begin{theorem}[Abel's Theorem: Second Form (Monotone Coefficients)]
If $a_n\ge 0$ and $\sum a_n$ diverges, then $\sum a_n x^n \to +\infty$ as $x\uparrow 1$.
(More generally, Abel gives boundary control for power series as $x\to 1^-$.)
\end{theorem}

\begin{example}[Deduce $\log 2 = 1-\frac12+\frac13-\frac14+\cdots$]
\end{example}

\begin{proof}[Solution]
For $|x|<1$, the geometric series gives
\[
\frac{1}{1+t} = \sum_{n=0}^\infty (-1)^n t^n.
\]
Integrate from $0$ to $x$:
\[
\int_0^x \frac{1}{1+t}\,dt = \sum_{n=0}^\infty (-1)^n \int_0^x t^n\,dt
\quad\Rightarrow\quad
\log(1+x)=\sum_{n=0}^\infty (-1)^n \frac{x^{n+1}}{n+1}.
\]
So for $|x|<1$,
\[
\log(1+x)= x-\frac{x^2}{2}+\frac{x^3}{3}-\frac{x^4}{4}+\cdots.
\]
Now apply Abel's theorem at $x=1$ (the alternating harmonic series converges), giving
\[
\log 2 = 1-\frac12+\frac13-\frac14+\cdots.
\]
\end{proof}

\begin{example}[Compute and prove the series for $\frac12(\log(1+x))^2$]
Show for $-1<x\le 1$:
\[
\frac12[\log(1+x)]^2
= \frac{x^2}{2}
- \frac{x^3}{3}\left(1+\frac12\right)
+ \frac{x^4}{4}\left(1+\frac12+\frac13\right)
-\cdots .
\]
\end{example}

\begin{proof}[Solution]
For $|x|<1$,
\[
\log(1+x)=\sum_{n=1}^\infty (-1)^{n-1}\frac{x^n}{n}.
\]
Square using the Cauchy product (valid for $|x|<1$ since the series is absolutely convergent for $|x|<1$? 
Here it is not absolutely convergent at $x=1$, but for $|x|<1$ it is absolutely convergent indeed because
$\sum |x|^n/n$ converges). Thus for $|x|<1$,
\[
(\log(1+x))^2
= \sum_{n=2}^\infty c_n x^n,
\quad
c_n=\sum_{k=1}^{n-1} (-1)^{k-1}\frac{1}{k}\cdot (-1)^{(n-k)-1}\frac{1}{n-k}.
\]
Combine signs:
\[
(-1)^{k-1}(-1)^{n-k-1}=(-1)^{n-2}=(-1)^n,
\]
so
\[
c_n = (-1)^n \sum_{k=1}^{n-1}\frac{1}{k(n-k)}.
\]
Now use the identity
\[
\frac{1}{k(n-k)}=\frac{1}{n}\left(\frac{1}{k}+\frac{1}{n-k}\right).
\]
Therefore
\[
\sum_{k=1}^{n-1}\frac{1}{k(n-k)}
=\frac{1}{n}\sum_{k=1}^{n-1}\left(\frac{1}{k}+\frac{1}{n-k}\right)
=\frac{1}{n}\cdot 2\sum_{k=1}^{n-1}\frac{1}{k}.
\]
Let $H_{n-1}:=\sum_{k=1}^{n-1}\frac{1}{k}$ (harmonic number). Then
\[
c_n = (-1)^n \cdot \frac{2}{n}H_{n-1}.
\]
Hence for $|x|<1$,
\[
(\log(1+x))^2 = \sum_{n=2}^\infty (-1)^n \frac{2H_{n-1}}{n} x^n,
\]
so
\[
\frac12(\log(1+x))^2 = \sum_{n=2}^\infty (-1)^n \frac{H_{n-1}}{n} x^n.
\]
Writing out terms:
\[
n=2:\quad (+)\frac{H_1}{2}x^2=\frac{1}{2}x^2,
\]
\[
n=3:\quad (-)\frac{H_2}{3}x^3=-\frac{1+\frac12}{3}x^3,
\]
\[
n=4:\quad (+)\frac{H_3}{4}x^4=\frac{1+\frac12+\frac13}{4}x^4,
\]
and so on, which matches the required expansion.

Finally, to extend from $|x|<1$ to $-1<x\le 1$, use Abel-type boundary control:
the power series defines a continuous function on $(-1,1]$ in this alternating setting,
and the identity holds on $(-1,1)$; taking $x\uparrow 1$ preserves convergence of the alternating sums,
yielding the formula for $-1<x\le 1$.
\end{proof}

%%%%%%%%%%%%%%%%%%%%%%%%%%%%%%%%%%%%%%%%%%%%%%%%%%%%%%%%%%%%
\end{document}
